% Worked solutions --- the full worked solution set for one assigned exercise set.
% Copy it into the meeting's 20 Exercises/<set>/ folder and replace every <...> placeholder.
\documentclass[a4paper,10pt]{article}
\IfFileExists{preamble.tex}{\input{preamble.tex}}{\IfFileExists{../../../../../60 Templates/preamble.tex}{\input{../../../../../60 Templates/preamble.tex}}{\IfFileExists{../../../../60 Templates/preamble.tex}{\input{../../../../60 Templates/preamble.tex}}{\input{../../../60 Templates/preamble.tex}}}}
\begin{document}
\ArtifactHeader{\ProjectCode}{<meeting and existing Research-map key or set>}{<the exercise set's name> --- Solutions}{<YYYY-MM-DD>}%
  {<the assigned source and any verified solutions checked against>}

\section{<The source's own grouping, or Solutions where it has none>}

% The optional argument takes the source's own numbering, which is what a reader looks for.
\begin{question}[<the source's number>]
<The question, restated closely enough that the solution reads without the source in hand.>
\end{question}

\begin{solution}
<The working, in the steps a reader would have to produce themselves --- the step that is easy to
skip is the step worth writing.>
\[
  \text{<the display the answer turns on>}
\]
<The answer, stated as an answer.>
\end{solution}

\begin{question}[<the source's number>]
<A question whose difficulty is a trap rather than the algebra.>
\end{question}

\begin{solution}
<The working.>
\end{solution}

\begin{warning}
<The trap that question sets, and how to see it coming next time. This belongs here only when it
is about the question; what the exercise set taught goes in the concepts consolidation.>
\end{warning}
\end{document}
